% Source: https://artofproblemsolving.com/community/c3h2221402_super_nice_problem
% Stable identity: aops-topic-2221402-q01
Since $AB=AC$ and $\angle A=120^\circ$, the base angles of triangle $ABC$ are $\angle ABC=\angle ACB=30^\circ$. By the congruence, $FC=FD=AB=AC$ and $\angle FCD=\angle ACB=30^\circ$.

\begin{asy}
usepackage("amsmath");
size(6cm);
defaultpen(fontsize(10pt));
real s=sqrt(6);
pair B=(0,0), C=(s,0), D=(s,s), Ept=(0,s);
pair A=(s/2,-s/(2*sqrt(3)));
pair F=(s-s/(2*sqrt(3)),s/2);
draw(B--C--D--Ept--cycle,linewidth(0.8));
filldraw(A--B--C--cycle,lightblue+opacity(0.4),blue);
filldraw(F--C--D--cycle,lightred+opacity(0.4),red);
draw(A--F,dashed);
pair u=unit(A-C), w=unit(F-C); draw(C+0.25*u--C+0.25*(u+w)--C+0.25*w);
dot(A);dot(B);dot(C);dot(D);dot(Ept);dot(F);
label("$B$",B,W);label("$C$",C,SE);label("$D$",D,NE);label("$E$",Ept,NW);label("$A$",A,S);label("$F$",F,W);
label("$120^\circ$",A+(0.33,0.33),blue+fontsize(8pt));label("$120^\circ$",F,E,red+fontsize(8pt));
label("$30^\circ$",C+0.75*dir(196),fontsize(7pt));label("$30^\circ$",C+0.75*dir(104),fontsize(7pt));
label("$2000$",(A+F)/2,NW);
\end{asy}

\textbf{Angle $ACF$.} The ray $CA$ makes $30^\circ$ with $CB$ on the outside of the square, and the ray $CF$ makes $30^\circ$ with $CD$ on the inside. Hence
\[\angle ACF=\angle ACB+\angle BCF=30^\circ+(90^\circ-30^\circ)=90^\circ.\]
So triangle $ACF$ is an isosceles right triangle with legs $AC=CF$ and hypotenuse $AF=2000$, and $AC=\dfrac{2000}{\sqrt2}=1000\sqrt2$.

\textbf{Side of the square.} Dropping the altitude from $A$ to $BC$ splits triangle $ABC$ into two $30$-$60$-$90$ triangles with hypotenuse $AC$, so $BC=2\cdot\dfrac{\sqrt3}{2}AC=\sqrt3\,AC=1000\sqrt6$.

The area of the square is $BC^2=\left(1000\sqrt6\right)^2=\boxed{6\,000\,000}$.
